\documentclass[11pt,a4paper]{article}

\usepackage[utf8]{inputenc}
\usepackage[T1]{fontenc}
\usepackage[margin=1in]{geometry}
\usepackage{graphicx}
\usepackage{booktabs}
\usepackage{amsmath,amssymb}
\usepackage{hyperref}
\usepackage{enumitem}
\usepackage{caption}
\usepackage{subcaption}
\usepackage{xcolor}
\usepackage{float}
\usepackage{array}
\usepackage{longtable}

\hypersetup{
  colorlinks=true,
  linkcolor=blue,
  citecolor=blue,
  urlcolor=blue
}

\title{Quantum Federated Learning with Quantum Fisher Information-Based Machine Unlearning on FEMNIST}
\author{Renato D. Souza}
\date{}

\begin{document}

\maketitle

\begin{abstract}
This paper presents a structured Quantum Federated Learning (QFL) framework built on PennyLane for a federated architecture composed of one central server and five clients. The study considers the FEMNIST dataset and evaluates two experimental settings: pure federated quantum training and machine unlearning using Quantum Fisher Information (QFI) after disconnecting one client from the federation. We define evaluation criteria for the unlearning process using forget-set accuracy, retain-set accuracy, membership inference attack success rate, and QFI trace. The manuscript provides a reproducible LaTeX skeleton aligned with the implementation artifacts of the project and is designed to be populated with empirical results from the final experimental pipeline.
\end{abstract}

\section{Introduction}

Quantum Federated Learning (QFL) extends the federated learning paradigm by replacing or augmenting classical client models with quantum or hybrid quantum-classical models \cite{fedavg,kairouz2019,qfl_towards_2023,qfl_foundations_2023}. In this work, we study QFL under a privacy-sensitive scenario based on FEMNIST, a benchmark designed for federated settings with heterogeneous and non-i.i.d.\ client data \cite{leaf}. The central objective is twofold: first, to train a global quantum model across five federated clients; second, to investigate machine unlearning when one client must be removed from the federation and its historical contribution should no longer influence the global model.

Machine unlearning is particularly relevant in privacy-preserving learning systems because data subjects or organizations may request removal from a trained model \cite{unlearning_survey_2024}. In a quantum setting, this poses additional challenges due to the structure of variational circuits and the geometry of the parameter space. We therefore adopt Quantum Fisher Information (QFI) as a signal for sensitivity and model influence during the unlearning phase \cite{qfi,federated_unlearning_active_forgetting,forgettable_federated_linear_learning}.

The project is implemented as a clean and reproducible software stack in Python 3.10+, with a central server connected to five clients, a PennyLane-based quantum model, and GPU-aware execution through `lightning.gpu` when available, falling back to `lightning.qubit` or `default.qubit` otherwise \cite{pennylane,lightning2024}. The codebase is organized into separate experiment folders for pure QFL training and QFI-based unlearning, with checkpoints, logs, outputs, and test coverage that support reruns and incremental development.

The experimental design follows two controlled settings. In the first, the model is trained for three federated rounds over five clients. In the second, the federation is trained and then one client, `client_0`, is excluded and the model is retrained over the remaining four clients for three rounds, after which the unlearning signal is summarized using `forget_set_accuracy`, `retain_set_accuracy`, `mia_success_rate`, and `qfi_trace`. The FEMNIST images are reduced to a compact four-feature representation via quadrant compression to make the circuit tractable while preserving the structure of the federated pipeline.

The paper is organized as follows. Section~\ref{sec:background} summarizes the background on federated learning, QFL, PennyLane, and QFI. Section~\ref{sec:methodology} introduces the architecture and algorithmic design. Section~\ref{sec:experiments} details the experimental protocol. Section~\ref{sec:results_expected} defines the expected outputs and evaluation metrics. Section~\ref{sec:discussion} discusses limitations and research directions.

\section{Background}
\label{sec:background}

Federated learning enables multiple clients to collaboratively train a model without centralizing raw data \cite{fedavg,kairouz2019}. In the FEMNIST benchmark, each client corresponds naturally to a user or device, creating a heterogeneous and non-i.i.d.\ data distribution. This setting is well suited for studying robust aggregation and privacy-aware learning \cite{leaf}.

PennyLane provides a practical framework for implementing variational quantum circuits with automatic differentiation \cite{pennylane}. In this project, a small-width quantum model is used to keep the circuit tractable while preserving the experimental structure required for federated training. The architecture uses a GPU-capable backend when available, falling back to CPU execution otherwise, which aligns with the Lightning simulator family used for accelerated hybrid quantum computation \cite{lightning2024}.

Quantum Fisher Information captures how sensitive a quantum state is to parameter changes \cite{qfi}. In the context of unlearning, we use QFI as a geometric proxy for the influence of training data and parameter updates. The hypothesis is that successful unlearning should reduce the dependence of the global model on the removed client's contribution while preserving performance on retained clients. This interpretation is consistent with recent work on machine unlearning, including federated unlearning strategies and certified unlearning formulations \cite{unlearning_survey_2024,federated_unlearning_active_forgetting,forgettable_federated_linear_learning}.

Membership inference attacks provide a practical privacy probe for evaluating whether a model still retains identifiable information about specific participants \cite{mia}. For this reason, the unlearning experiment combines accuracy-based and privacy-based metrics rather than relying on a single geometric quantity.

\section{Methodology}
\label{sec:methodology}

The proposed system consists of one central server and five clients. The server coordinates rounds of training, while each client performs local updates using a PennyLane-based quantum classifier. After local computation, the server aggregates updates using a sample-size-weighted strategy, following the standard federated optimization principle introduced in FedAvg \cite{fedavg}.

The experimental pipeline is divided into two parts. In the first experiment, the system performs pure QFL training over all five clients. In the second experiment, one client is excluded from the federation and the model is retrained on the remaining clients. The unlearning procedure uses QFI to quantify the sensitivity of the updated model and supports evaluation of whether the disconnected client has been effectively forgotten \cite{qfi,federated_unlearning_active_forgetting,forgettable_federated_linear_learning}.

The FEMNIST data are reduced to a compact four-feature representation to make the quantum circuit feasible. Although this representation is a simplification of the original image structure, it allows controlled experimentation with the federated and unlearning workflows while the full ingestion pipeline is developed \cite{leaf}. This reduction is implemented as quadrant compression over normalized 28 \(\times\) 28 images.

The implementation is intentionally structured for reproducibility. Both experiments use `seed = 42`, `num_clients = 5`, and `num_rounds = 3`. The unlearning experiment excludes `client_0`, yielding four active clients in the post-unlearning retraining phase. The scripts read the intermediate dataset from `data/femnist_sample.npz`, write JSON summaries to dedicated output directories, and maintain checkpoints so that interrupted runs can be resumed without discarding previous progress.

From a software-engineering perspective, the repository is organized around a clean separation of concerns: dataset utilities live in `src/qfl/data`, federated orchestration in `src/qfl/federated`, quantum circuit logic in `src/qfl/quantum`, and utility code such as deterministic seeding in `src/qfl/utils`. This structure keeps the experiment scripts concise and makes the project easier to test and extend.

\section{Experimental Design}
\label{sec:experiments}

Two experiments are considered, both using the same dataset preparation, client partitioning strategy, and deterministic seed.

\subsection{Experiment I: Pure QFL Training}

The first experiment trains a global quantum model using the full federation of five clients. The input data are loaded from `data/femnist_sample.npz`, normalized, compressed into four quadrant features, and partitioned into `num_clients = 5` client splits. The model is trained for `num_rounds = 3` with `seed = 42`. This setting evaluates convergence behavior, aggregation stability, and the feasibility of quantum training under heterogeneous client partitions. The use of a compact variational circuit is motivated by the current state of PennyLane-based QML tooling and the need to keep the circuit tractable during iterative federated updates \cite{pennylane,lightning2024,qfl_towards_2023}.

\subsection{Experiment II: QFI-Based Machine Unlearning}

The second experiment disconnects one client from the federation and re-evaluates the global model after retraining on the remaining clients. In the current implementation, `client_0` is excluded, leaving four active clients. The unlearning procedure is examined using `forget_set_accuracy`, `retain_set_accuracy`, `membership inference attack success rate`, and `qfi_trace`. The training and unlearning phases are both executed for `num_rounds = 3` with `seed = 42`. The design follows the literature on federated unlearning and certified data removal, but adapts it to a quantum setting where QFI provides a geometry-aware proxy for model influence \cite{unlearning_survey_2024,federated_unlearning_active_forgetting,forgettable_federated_linear_learning,qfi}.

The evaluation is designed to answer the following questions:
Does the model lose predictive power on the forgotten client's data? Does the model preserve performance on the retained clients? Does unlearning reduce exposure to membership inference attacks? Is there a measurable reduction in the influence of the removed client according to QFI? These questions align with the privacy and completeness criteria commonly used in machine unlearning studies \cite{mia,unlearning_survey_2024}.

The implementation writes the training summary to `experiments/qfl_training/outputs/training_summary.json` and the unlearning summary to `experiments/qfl_unlearning_qfi/outputs/unlearning_summary.json`. Checkpoints are also maintained for both experiments so that execution can be resumed from the last completed round.

\section{Expected Results}
\label{sec:results_expected}

This manuscript is currently structured around the expected outcomes of the experiments rather than final empirical values. The tables below define the result template that will be populated once the full FEMNIST pipeline is executed. Because the current codebase is designed for reproducibility, the expected outputs are already tied to fixed experimental parameters: `seed = 42`, `num_clients = 5`, `num_rounds = 3`, and `excluded_client_id = client_0` in the unlearning setting.

\subsection{Expected Training Summary}

\begin{table}[H]
\centering
\caption{Expected summary of the pure QFL training experiment.}
\label{tab:training_expected}
\begin{tabular}{lccc}
\toprule
Metric & Round 1 & Round 2 & Round 3 \\
\midrule
Number of active clients & 5 & 5 & 5 \\
Aggregation stability & High & High & High \\
Global model availability & Yes & Yes & Yes \\
GPU backend usage & Preferred & Preferred & Preferred \\
\bottomrule
\end{tabular}
\end{table}

\subsection{Expected Unlearning Summary}

\begin{table}[H]
\centering
\caption{Expected metrics for the QFI-based unlearning experiment.}
\label{tab:unlearning_expected}
\begin{tabular}{lcc}
\toprule
Metric & Target behavior & Interpretation \\
\midrule
Forget-set accuracy & Decrease toward random baseline & Removed client is effectively forgotten \\
Retain-set accuracy & Remain stable or slightly decrease & No catastrophic forgetting \\
MIA success rate & Decrease & Reduced privacy leakage \\
QFI trace & Reduce after unlearning & Lower sensitivity to removed client \\
\bottomrule
\end{tabular}
\end{table}

\subsection{Placeholder Result Schema}

The following JSON-like schema reflects the expected output generated by the unlearning script:

\begin{verbatim}
{
  "qfi_trace": <float>,
  "remaining_clients": 4.0,
  "excluded_client": "client_0",
  "forget_set_accuracy": <float>,
  "retain_set_accuracy": <float>,
  "random_baseline_accuracy": <float>,
  "mia_success_rate": <float>
}
\end{verbatim}

The final paper should replace placeholders with actual values, confidence intervals, and, if available, multiple runs across random seeds.

When the experiments are executed on the RunPod environment described in the repository, the resulting summaries should be cross-checked against the checkpoint files and logs to verify that the same deterministic configuration was used end-to-end.

\section{Discussion}
\label{sec:discussion}

The current implementation prioritizes a reproducible and modular architecture over immediate benchmark-level performance. This choice is appropriate for an initial research scaffold, but several limitations remain. First, the FEMNIST ingestion pipeline is simplified through feature compression, so the model does not yet exploit the full spatial structure of the original images \cite{leaf}. Second, the quantum model is intentionally compact and may not capture the expressivity required for high-accuracy classification. Third, the unlearning metric based on QFI should be interpreted as a geometric indicator rather than a complete proof of erasure \cite{qfi,unlearning_survey_2024}.

Despite these limitations, the project establishes a clear experimental path: it separates training and unlearning, integrates GPU-aware execution in PennyLane, and introduces a privacy-oriented evaluation protocol based on forget-set performance, retain-set performance, membership inference risk, and QFI sensitivity. From the perspective of research contribution, the main value of the current repository lies in providing a clean software scaffold that can be extended toward stronger quantum ansätze, richer federated aggregation strategies, and more formal unlearning guarantees.

A further open problem is the lack of a full FEMNIST ingestion pipeline directly from LEAF. Until that is implemented, the current `npz` intermediary should be treated as an experimental convenience rather than as a final benchmark protocol. Future work should also validate whether the QFI-based signal correlates with stronger empirical unlearning criteria under multiple random seeds and client-dropout scenarios \cite{federated_unlearning_active_forgetting,forgettable_federated_linear_learning,qfl_foundations_2023}.

\section{Conclusion}
\label{sec:conclusion}

This paper documents a QFL research pipeline with a dedicated machine unlearning phase using QFI. The LaTeX structure is aligned with the implementation artifacts in the repository and is intended to be populated by empirical results from the two controlled experiments described above. The main contribution of the current repository is not a final performance claim but a clean, reproducible framework for studying quantum federated training and unlearning on FEMNIST \cite{leaf,qfl_towards_2023,qfl_foundations_2023}.

The present version already fixes the core experimental parameters, including a five-client federation, three training rounds, a deterministic seed of 42, and a designated forgotten client (`client_0`). This makes the project suitable as a research baseline that can later be extended with a full LEAF ingestion pipeline, more expressive quantum circuits, and more rigorous unlearning evaluations.


\bibliographystyle{plain}
\bibliography{references}

\end{document}

